\documentclass{jsarticle}
\usepackage[dvipdfmx,final]{graphicx}
\usepackage{chemfig}
\usepackage{amsmath}
\usepackage{amssymb}
\begin{document}
\title{Time measure of global scale revearsed with\\
space measure of variable metric}
\author{Masaaki Yamaguchi}
\date{}
\maketitle
\newcommand{\variint}{%
	\begin{picture}(15,30)
		\put(0,0){
			$ \iint $}
		\put(0,5){\line(15,0){15}}
	\end{picture} %
}
   \newcommand{\squareiint}{%
	\begin{picture}(15,30)
		\put(0,0){
			$ \iint $}
		\put(0,5){$\text{\scalebox{2.0}[1]{$\square$}}$}
	\end{picture} %
}

  
   \newcommand{\dazanier}{%
	\begin{picture}(15,30)
		\put(0,0){
			$ \bigsqcup $}
		\put(4,3){$\text{\scalebox{1.0}[1]{$\top$}}$}
		\put(4,3){$\text{\scalebox{1.0}[1]{$\bot$}}$}
	\end{picture} %
}


\newcommand{\variiint}{%
	\begin{picture}(15,15)
		\put(0,0){
			$ \iiint $}
		\put(0,5){\line(15,0){20}}
	\end{picture} %
}



 \newcommand{\alearletter}{%
	 \begin{picture}(10,20)
		 \put(0,0){
			 $ \square $}
		 \put(0,0){\line(1,1){10}}
	 \end{picture} %
	 }

  \newcommand{\whitewithbodercircle}
  {\text{\textcircled{\phantom{\textbullet}}}
  \kern-0.9em\text{\textcircled{\phantom{}}}}
     

 
  \newcommand{\whitewithblack}{\text{\textcircled{\textbullet}}}
  \newcommand{\whitewithblacksquare}{\text{\fbox{\text{\rule{0.5em}
  {0.5em}}}}}
  \newcommand{\whitewithwhitesquare}{\text{\fbox{\text{\phantom
  {\rule{0.5em}{0.5em}}}}}}
 \newcommand{\whitewithbodercircleiiint}{%
	 \begin{picture}(30,45)
		 \put(0,0){\scalebox{2.0}[2]{$\whitewithbodercircle$}}
		 \put(0,0){\scalebox{2.0}[2]{$\iiint$}}
	 \end{picture}%
	 }
\newcommand{\whitewithbodercircleint}{%
	 \begin{picture}(30,45)
		 \put(0,0){\scalebox{2.0}[2]{$\whitewithbodercircle$}}
		 \put(0,0){\scalebox{2.0}[2]{$\int$}}
	 \end{picture}%
	 }
 \newcommand{\whitewithbodercirclevarint}{%
	 \begin{picture}(30,45)
		 \put(0,0){\scalebox{2.0}[2]{$\whitewithbodercircle$}}
		 \put(0,0){\scalebox{2.0}[2]{$\varint$}}
	 \end{picture}%
	 }
 
 \newcommand{\whitewithbodercirclevariiint}{%
	 \begin{picture}(30,45)
		 \put(0,0){\scalebox{2.0}[2]{$\whitewithbodercircle$}}
		 \put(0,0){\scalebox{2.0}[2]{$\variiint$}}
	 \end{picture}%
	 }
 
 \newcommand{\average}{%
	  \begin{picture}(30,45)
		  \put(0,0){\scalebox{2.0}[2]{$\square$}}
		  \put(4,4){$\setminus$}
	  \end{picture}%
  } 
 
 These equation are Projection function, cross time, assemble of projection
 function, estimate manifold, global integral manifold, project global 
 integral manifold, average of add and even function.
 This system is signal of write with system.
  $$ \whitewithblack f_m , f \boxtimes g = F^{m}_t, 
  \whitewithblacksquare f_m = {\begin{pmatrix}
	  n \\
  r \end{pmatrix}}^{\otimes L} = \lim_{n \to \infty}
  {\begin{pmatrix}
	  n \\
  r \end{pmatrix}}^{\otimes L} = {d \over {d f}}\boxtimes $$

  \newcommand{\whitewithwhitecircle}
  {\text{\textcircled{\phantom{\textbullet}}}}
  Estimate function is equal with dervargence function
  $$ \begin{pmatrix}
	  n \\
  r \end{pmatrix}(\zeta(s)) = \whitewithbodercircle f_m $$
  $$ \whitewithblack f_m ,
  \whitewithblacksquare f_m = {\begin{pmatrix}
	  n \\
  r \end{pmatrix}}^{\otimes L} = \lim_{n \to \infty}
  {\begin{pmatrix}
	  n \\
  r \end{pmatrix}}^{\otimes L} = {d \over {d f}}\boxtimes $$


 

$$ \whitewithbodercircleint \sum [dI_m] $$
$$ {}^t\variiint_{D\chi} cohom D\chi = \log(x\log x) = [x_m,y_m]\times
[x_n,y_n] $$
$$ \int R^{\nabla r}dr_x = l^{R|\nabla} = \beta(p,q) $$
$$ V_{D\chi}\int dx_m = S \otimes S = \vee\int dx_m = S^2 \otimes S^2 $$
$$ \Delta[\vartriangleleft_x,\vartriangleleft_y] = \average(x,y) $$

  \newcommand{\whitewithboder}{%
	  \begin{picture}(15,30)
		  \put(0,0){\scalebox{2.0}[2]{$\square$}}
		  \put(4,4){$\square$}
	  \end{picture}%
  } 
  \newcommand{\whitewithnabla}{%
	  \begin{picture}(15,30)
		  \put(0,0){\scalebox{2.0}[2]{$\nabla$}}
		  \put(3,3){$\nabla$}
	  \end{picture}%
  } 


  Assemble function is Euler equation, and this equation also
  Euler product of fundermantal manifold.
  $$ \whitewithnabla = e^{i\theta} = \cos \theta + i \sin \theta $$

  Average function.
  $$ \average = {{f(x) + g(x)}} \ge 2{\sqrt{f(x)g(x)}} $$
  Global manifold is partial integral of global topology.
 $$ \whitewithblacksquare f_m =  \whitewithboder \text{                 
      } f_m = 
 \whitewithbodercircle f_m = 
  D_{\mathcal{K}}(x) = \int dx_m , \mathrm{ker}/\mathrm{im}f = 
  \partial $$

$$\boxtimes$$
$$\whitewithbodercircle, \whitewithblack $$
$$\whitewithblacksquare,\whitewithwhitecircle$$
$$\scalebox{2.0}[2]{\variint}$$ 
$$\whitewithboder$$ 
$$ \whitewithnabla$$
$$\average$$ 
は、それぞれ、積のクロス、 大域的積分多様体、組み合わせ多様体、
同型の大域的積分多様体、準同型の大域的積分多様体、
準同型写像、大域的切断、重ね合わせの原理、
相加相乗平均、をそれぞれ表す。
 

$$||ds^2||=\whitewithbodercircleiiint \mathrm{Docol}[I_m]$$
$$  \whitewithbodercircleint \sum[dI_m]$$

$$ {}^t\scalebox{2.0}[2]{\variint}_{D\chi}\mathrm{colm}D\chi = \log(x\log x)$$
$$ = [x_m,y_m]\times [x_n,y_n] $$

\newcommand{\timebow}{%
	\begin{picture}(5,5)
		\put(0,0){$\text{\rotatebox{90}{$\bowtie$}}$}
	\end{picture} %
	}


 \newcommand{\flowerletter}{%
	 \begin{picture}(10,20)
		 \put(0,0){$\text{\scalebox{1.0}[1]{$\square$}}$}
		 \put(0,0){$\text{\scalebox{1.0}[2]{$\int$}}$}
	 \end{picture} %
	 }



       時間計量による、指数定理における擬微分と擬積分多様体を微小変量によって、
   大域的トポロジーが成されられる。ガンマ関数における大域的多様体と、
   ヒッグス方程式における大域的多様体の計算式を以下に、載せておく。
   $$ {T^{\mu\nu}}^{T^{\mu\nu}} = \int T^{\mu\nu} \rotatebox{90}{$\bowtie$}_m $$
   $$   {T^{\mu\nu}}^{{T^{\mu\nu}}^{'}}=\int T^{\mu\nu}dx_m  $$
      $$ {\Gamma^{\gamma}}^{'} = {d \over {d \gamma}}\Gamma $$
   $$ \Gamma^{\gamma} = \int \Gamma dx_m $$
   $$ \rotatebox{90}{$\bowtie$} = \int x^x dx_m $$


     ヒッグス場方程式が、相加相乗平均方程式であり、大域的多様体でもある。
  $$ {d \over {df}}F(x,y)=m(x) $$
  $$ = {d \over {d m(x)}}M = \partial (\sigma(M)) $$
  $$=  {d \over {d \timebow}}Z^{'}(\zeta)dx_m $$
  $$ ||ds^2|| = \nabla_{i}\nabla_{j}\int \nabla M(\sigma)d\eta $$
  M理論の指数定理による擬微分と擬積分多様体が、ヘルマンダー作用素から
  のノルム空間となり、
  $$ = M^{{m(x)^{'}}} $$
  M理論のニュートン方式の大域的微分多様体になる。
  $$ ||ds^2|| = \int_M [d(\partial) + d(\sigma)]dI_m $$
  それが、D-braneのM理論へと導かれる。
  $$ {d \over {df}}(e^{-f} + e^{f}) = {{d \over {d R^{'}}}}R = d[I_m] $$
  結局は、Jones多項式における、大域的多様体としての、虚数微分方程式となる。
  リッチ・フロー方程式による、大域的多様体の擬微分と擬積分多様体の
  計算式が、ニュートン方式による、擬積分と擬微分多様体となり、
  時間計量が、相対性による、測量がなくなり、空間計量となり、
  空間概念の量子化となる。
$$ ||ds^2|| = \whitewithbodercircleiiint Docol[I_m] $$

 	
$$ \timebow = \flowerletter r dr_m $$
$$ \dazanier = D_{\chi}{}^{[t}\variiint \mathrm{exp}[\mathrm{cohom}
D_{\chi}[I_m] $$ 
$$ D_{\chi} |: x \to y $$
$$ D_{\chi}(\chi)y = \log D_{\chi}^{y} = x $$
$$ \square \cap \Psi \ni  \zeta(s) $$
$$ \int \kappa (A^{\mu\nu})^{2}dx_m = \int \Gamma^{'}(\gamma)dx_m $$
$$ g_{\mu\nu}(x)g^{-1}_{\mu\nu} = {}^{t}\variiint_{D\chi}\kappa
(A^{\mu\nu})dx_m $$
$$  \mathcal{S}^{\mu\nu}_{D(\chi,D)} = \bigoplus {(V_{k}{}^{N}\squareiint
\Gamma(\gamma)d\gamma)}[I_m] $$
$$ \mathcal{O}(x) = \oint \pi x^2 dx $$
$$ \int e^{-x}x^{1-t}dx = \int \sin \theta \cos \theta d \theta \int e^{\sin \theta \cos \theta}d\theta $$


\end{document}
